\chapter{Referenz: Implementierung, Datentafeln, Registerstand}
\label{ch:anhang-referenz}

\section*{Die Referenzimplementierung}

Alle Enumerationen und Verifikationen dieses Buchs beruhen auf
einer kleinen, bewusst unoptimierten Python-Implementierung
(\texttt{horimetrik\_reference.py}) mit Selbsttests: Codierung und
Decodierung, Struktur\-polynome, beide Erweiterungen, beide
Kompaktifizierungen, Strukturarithmetik. Ergänzt wird sie durch
datierte Experimentskripte, deren Ergebnisse im Arbeitsregister mit
Status (bewiesen / vermutet / widerlegt) geführt werden. Der
Grundsatz des Projekts: Enumeration erzeugt Vermutungen und
Gegenbeispiele, niemals Sätze -- und jede behauptete Schranke muss
konstruierte Extremfälle überleben, nicht nur Stichproben.

\section*{Die neun Zählfolgen}

\begin{center}
\footnotesize
\begin{tabular}{lll}
\toprule
Folge & Anfangswerte & Status\\
\midrule
Sprungstellen $m_g$ ($g{\ge}1$) & $3,15,84,533,3883,31998,294875,\ldots$ &
Tiefengesetz, Asymptotik bewiesen\\
$\beta(r,r)$ ($r{\ge}1$) & $1,6,22,87,407,2473,19527,\ldots$ & Asymptotik bewiesen\\
Springer $D(n)$ ($n{\ge}3$) & $1,8,60,360,2520,21168,211680,2268000$ &
Ziffernformel bewiesen\\
Nicht-Springer $N(n)$ ($n{\ge}3$) & $5,16,60,360,2520,19152,151200,\ldots$ &
Ziffernformel bewiesen\\
Fixpunkte $F(n)$ ($n{\ge}2$) & $4,17,88,540,3960,32760$ & $=n\cdot n!-D(n)$ (bewiesen)\\
Ein-Ziffern-Fakultäten & $1,2,5,7,11,23,29,39,59,119,\ldots$ &
vollständig charakterisiert\\
$\max\Omega$ ($n{\ge}1$) & $0,1,6,43,351,2873,26443,270291$ & $\Omega\ge0$
bewiesen\\
$\#\{\Omega=0\}$ ($n{\ge}1$) & $2,5,13,23,61,111,274,564$ &
Gleichheitskriterium bewiesen\\
$\Gamma_\circ(r,r)$ ($r{\ge}1$) & $1,4,23,6541,477149751,\ldots$ &
Extremalprinzip bewiesen\\
\bottomrule
\end{tabular}
\end{center}

Keine dieser Folgen war zum Zeitpunkt der Niederschrift in der
OEIS; die Einreichung ist vorbereitet (A-Nummern werden hier
nachgetragen).

\section*{Bestätigte Vorhersagen}

\begin{enumerate}
\item Tiefe $-3$ des Kommutators erstmals nahe Schale $100$
      (exakt: 84 als Sprungstelle, bestätigt im Bereich bis 200).
\item $m_6=31998$ bei prognostizierten $\approx31930$.
\item $m_7=294875$ -- exakt wie vorhergesagt.
\item Ein-Ziffern-Fakultäten bei $n=29,39,59,119$ und
      $143,179,239,359,719$ aus der (inzwischen vollständigen)
      Charakterisierung $n=(i{+}1)!/d-1$ -- alle neun bestätigt.
\end{enumerate}

\section*{Registerstand bei Redaktionsschluss dieser Fassung}

44 bewiesene Sätze -- darunter die Grenzraum\-sätze der drei
Treppen, die zwei dualen Halbringe, der Dualitäts\-satz und der
Seitentreue-Satz, das Kompaktifizierungs\-monoid, das
Tiefengesetz, die Ziffernformel und die Interchange-Positivität
--, dazu 9 dokumentierte Sackgassen, 3 widerlegte eigene
Vermutungen, 9 unbekannte Zählfolgen und 4 bestätigte
Vorhersage-Gruppen (12 Einzeltreffer, Liste oben). Die jeweils aktuelle Fassung von Register,
Vermutungen, Gegenbeispielen und Changelog liegt den Arbeitsdateien
des Projekts bei.
